% Elementary recovery sequence. The more general theorem and proof details
% belong in the backup slides, rather than this first explanation.
\section{Finding hidden groups}

\begin{frame}{A graph is visible; its group labels are missing}
\begin{center}
\Afig[height=3.9cm]{recovery-uncolored}
\end{center}
The six numbers name the vertices. We see every edge, but no group colors.
\begin{block}{The question}
Can the connections tell us which vertices belong together?
\end{block}
\end{frame}

\begin{frame}{Start with one very concrete neighbor rule}
\begin{columns}[c]
\column{.53\textwidth}
\Afig[height=3.3cm,width=\linewidth,keepaspectratio]{recovery-colored}
\column{.44\textwidth}
Every vertex has:
\begin{itemize}
\item \textbf{two} neighbors in its own group;
\item \textbf{one} neighbor in the other group.
\end{itemize}
These are exact counts, not averages.
\end{columns}
\source{This six-vertex graph illustrates the rule. The theorem later concerns growing random graphs.}
\end{frame}

\begin{frame}{Check the rule at a single vertex}
\begin{columns}[c]
\column{.56\textwidth}
\Afig[height=3.3cm,width=\linewidth,keepaspectratio]{recovery-local-count}
\column{.41\textwidth}
Look at blue vertex \textbf{2}.
\begin{itemize}
\item Vertices 1 and 3 are blue.
\item Vertex 5 is orange.
\end{itemize}
So its counts are
\[
       (\text{same},\text{across})=(2,1).
\]
\end{columns}
A proposed grouping is valid only if \emph{every} vertex passes its rule.
\end{frame}

\begin{frame}{A plausible coloring can fail the rule}
\begin{columns}[c]
\column{.56\textwidth}
\Afig[height=3.3cm,width=\linewidth,keepaspectratio]{recovery-wrong}
\column{.41\textwidth}
In this trial coloring, vertex \textbf{2} has:
\begin{itemize}
\item one neighbor of its own color;
\item two neighbors of the other color.
\end{itemize}
\begin{alertblock}{Reject this trial}
The required counts were $(2,1)$, not $(1,2)$.
\end{alertblock}
\end{columns}
Checking the ordinary neighbor counts gives a definite answer.
\end{frame}

\begin{frame}{The algorithm also tries different count rules}
For a degree-three graph, try tables such as these.
Rows name your group; columns name a neighbor's group.
\begin{columns}[t]
\column{.30\textwidth}
\begin{block}{One group}
\[
\begin{pmatrix}3\end{pmatrix}
\]
Three neighbors inside.
\end{block}
\column{.32\textwidth}
\begin{block}{Two groups}
\[
\begin{pmatrix}2&1\\1&2\end{pmatrix}
\]
Two inside; one across.
\end{block}
\column{.32\textwidth}
\begin{block}{Another rule}
\[
\begin{pmatrix}1&2\\2&1\end{pmatrix}
\]
One inside; two across.
\end{block}
\end{columns}
A fixed upper bound on groups and neighbors leaves only finitely many tables.
The true table is \emph{not} supplied to the algorithm.
\end{frame}

\begin{frame}{The recovery algorithm: five steps}
\begin{center}
\Afig[width=\textwidth]{recovery-workflow}
\end{center}
The hard part is organizing all possible assignments of vertices to groups.

The next slides show how the middle steps work for the specific
\textbf{two-inside, one-across} rule on our six-vertex graph.
\end{frame}

\begin{frame}{Turn a grouping into simple equations}
\begin{columns}[c]
\column{.52\textwidth}
\Afig[height=3.3cm,width=\linewidth,keepaspectratio]{recovery-marks}
\column{.45\textwidth}
Mark group A by $1$ and group B by $0$.
Call the mark at vertex $i$ \,$x_i$.

Adding the neighbors' marks counts neighbors in group A.

At vertex 2 this gives
\[
       x_1+x_3+x_5=1+x_2.
\]
\end{columns}
If $x_2=1$, the count must be 2. If $x_2=0$, it must be 1.
Write one such equation for each vertex and solve them together.
\end{frame}

\begin{frame}{First solve an easier arithmetic version}
\begin{columns}[c]
\column{.53\textwidth}
\Afig[width=\linewidth]{recovery-clock}
\column{.44\textwidth}
Use a three-position clock: after 2 comes 0.

Solve the equations with this arithmetic.
Linear algebra gives a compact description of \emph{all} their solutions.
\end{columns}
Every genuine $0/1$ grouping survives this step, but extra answers can appear.
For example, 1 and 7 look the same on this clock, though they are different integers.
\end{frame}

\begin{frame}{For this graph and this rule: nine clock solutions}
\begin{center}
\Afig[height=3.2cm,width=.96\textwidth,keepaspectratio]{recovery-nine-candidates}
\end{center}
Each boxed string gives the six proposed marks $x_1,\ldots,x_6$.

Seven answers contain a 2, so they are not $0/1$ groupings.
The two remaining answers must still pass the ordinary neighbor-count check.
\source{An exact enumeration of all $3^6$ markings verifies this finite example. It is not the asymptotic theorem.}
\end{frame}

\begin{frame}{Nine candidates become one grouping}
\begin{center}
\Afig[height=3.1cm,width=.92\textwidth,keepaspectratio]{recovery-same-partition}
\end{center}
Both $111000$ and $000111$ pass every ordinary neighbor-count check.
They differ only by exchanging the two group names.
\[
 \boxed{9\text{ clock answers}}\ \longrightarrow\
 \boxed{2\text{ valid colorings}}\ \longrightarrow\
 \boxed{1\text{ division into groups}}.
\]
\source{This count is for the specified $(2,1)$ rule on this particular graph. Other count rules are a separate search.}
\end{frame}

\begin{frame}{Which valid answer should be kept?}
\begin{center}
\Afig[height=3.3cm,width=.96\textwidth,keepaspectratio]{recovery-coarse-fine}
\end{center}
Both pictured answers have a consistent neighbor-count table.
The two-group answer keeps more information than the one-group answer.

The algorithm keeps a valid answer with the most groups, subject to the fixed bound $K$.
\source{These are two illustrative answers; this tiny graph can have other valid groupings. The theorem controls all alternatives in the growing random model.}
\end{frame}

\begin{frame}{Which random graphs does the main case cover?}
\begin{block}{A precise setting, in words}
Every vertex has the same fixed total degree $d\ge3$ and follows a hidden
neighbor-count table with at most $K$ groups.
\end{block}
\begin{itemize}
\item The table connects the groups; their positive size ratios stay fixed.
\item For each group or pair of groups, choose the links uniformly from
all simple graphs with the required counts.
\item These choices are independent.
\end{itemize}
Our example uses random degree-two links inside each group and a random
perfect matching across. Random lifts of fixed regular bases are also covered.
\end{frame}

\begin{frame}{The theorem makes the final choice reliable}
\begin{theorem}[Recovery, stated informally]
In the preceding random model, as the graph grows, with probability tending
to one every valid grouping into at most $K$ groups only merges whole true groups.
\end{theorem}
Therefore the most detailed valid grouping is the hidden one.

The algorithm is given the graph and fixed bounds on groups and degree.
It is \emph{not} given the actual group count, sizes, or neighbor-count table.
\begin{block}{Two separate guarantees}
Every answer on the list passes its exact rules.
Identifying which answer was originally hidden is the high-probability theorem.
\end{block}
\source{The precise statement and complete proofs are in the report's main recovery theorem and proof sections.}
\end{frame}

\begin{frame}{Why is the method computationally useful?}
The proof controls the \emph{average} number of intermediate candidates,
including rare graphs with many of them.

It bounds the expected total work by nearly $n^2$, where $n$ is the number of vertices.
\begin{center}
\Afig[height=2.5cm,width=.75\textwidth,keepaspectratio]{recovery-work-bound}
\end{center}
Doubling $n$ roughly multiplies the bound by four, apart from logarithmic factors.
This is an expected bound, not a speed benchmark or a bound for every input.
\end{frame}

\begin{frame}{What is the new idea, and where does it stop?}
\begin{block}{The useful combination}
Exact neighbor counts give equations. A complete search organized by those
equations can recover unknown groups without needing an isolated signal eigenvalue.
\end{block}
The result covers a broad family of sparse random graphs with fixed parameters,
including every fixed connected simple regular base of degree at least three.
\begin{itemize}
\item The algorithm also learns the count table and group sizes.
\item Exact verification rules out false accepted groupings.
\item Missing or added edges can break the equations: there is no noise theorem here.
\end{itemize}
The formal report proves the model-specific guarantee and explains its predecessors.
\end{frame}
